\subsection{Linearization and Local LTI Model}
\label{sec:linearization}

To linearize the model, the nominal straight-line operating condition, where both rovers are moving in the same direction at the same speed while maintaining a constant separation, is considered with:
\[
\rho=d^\ast,\qquad
\beta=\varepsilon=0,\qquad
v_f=v_t=v_0>0,\qquad
\omega_f=\omega_t=0,
\]
where $d^\ast$ is the desired minimal separation between both vehicles and $v_0$ is the nominal forward speed.

The state vector, control input, disturbance from the target's motion, and measured output are then defined as
\[
\mathbf{x}=
\begin{bmatrix}
\delta\rho & \beta & \varepsilon
\end{bmatrix}^{\mathsf T},
\qquad
\mathbf{u}=
\begin{bmatrix}
\delta v_f & \omega_f
\end{bmatrix}^{\mathsf T},
\qquad
\mathbf{w}=
\begin{bmatrix}
\delta v_t & \omega_t
\end{bmatrix}^{\mathsf T},
\]
and
\[
\mathbf{y}=
\begin{bmatrix}
\hat{\rho}-d^\ast\\
\hat{\beta}
\end{bmatrix},
\]
where $\delta\rho=\rho-d^\ast$, $\delta v_f=v_f-v_0$, and $\delta v_t=v_t-v_0$. The relative heading is $\varepsilon=\theta_t-\theta_f$, and the bearing is $\beta=\lambda-\theta_f$. 

The measured outputs are the estimated range deviation $\hat{\delta\rho}$ and relative bearing $\hat{\beta}$ expressed relative to the follower's forward direction.

The control inputs are virtual commands: $\delta v_f$ represents a deviation from the nominal follower speed, while $\omega_f$ represents the commanded follower yaw rate. The target-motion variables $\delta v_t$ and $\omega_t$ are treated as disturbances because the follower does not directly know the target's real-time speed and yaw rate. \autoref{tab:model_variables} showcases the state space variables, their physical interpretation, and associated units. 

Linearizing the nonlinear relative-kinematics model about the nominal operating condition gives
\begin{equation}
\dot{\mathbf{x}}
=\mathbf{A}\mathbf{x}
+\mathbf{B}\mathbf{u}
+\mathbf{E}\mathbf{w},
\qquad
\mathbf{y}
=\mathbf{C}\mathbf{x}+\mathbf{n},
\label{eq:lin}
\end{equation}
where $\mathbf{n}$ represents measurement-estimation errors (to be determined experimentally). The measurement equation applies when a valid acoustic estimate becomes available; measurements are not assumed to arrive at every controller update.

For the constant-speed operating condition, the local model is LTI:
\begin{equation}
\begin{aligned}
\mathbf{A}
&=
\begin{bmatrix}
0&0&0\\
0&0&a\\
0&0&0
\end{bmatrix},
&
\mathbf{B}
&=
\begin{bmatrix}
-1&0\\
0&-1\\
0&-1
\end{bmatrix},\\[6pt]
\mathbf{E}
&=
\begin{bmatrix}
1&0\\
0&0\\
0&1
\end{bmatrix},
&
\mathbf{C}
&=
\begin{bmatrix}
1&0&0\\
0&1&0
\end{bmatrix},
\end{aligned}
\label{eq:mats}
\end{equation}
with
\[
a=\frac{v_0}{d^\ast}.
\]


Refer to Appendix \ref{app:der} for the supplementary derivations and to \autoref{fig:ctrl} for the block diagram. 

\begin{table*}[!t]
\centering
\caption{State-space variables and their physical interpretation.}
\label{tab:model_variables}
\small
\renewcommand{\arraystretch}{1.15}
\begin{tabular}{p{0.15\textwidth} p{0.10\textwidth} p{0.55\textwidth} p{0.10\textwidth}}
\hline
\textbf{Category} & \textbf{Variable} & \textbf{Physical meaning} & \textbf{Unit}\\
\hline
State
& $\delta\rho$
& Deviation of the actual follower-target distance  from the desired separation $d^\ast$
& m\\

State
& $\beta$
& Relative bearing of the target measured from the follower's forward direction
& rad\\

State
& $\varepsilon$
& Relative heading between the target and follower
& rad\\

Control input
& $\delta v_f$
& Deviation of the follower's forward speed from the nominal speed $v_0$
& m/s\\

Control input
& $\omega_f$
& Commanded yaw rate of the follower
& rad/s\\

Disturbance
& $\delta v_t$
& Deviation of the target's forward speed from the nominal speed $v_0$
& m/s\\

Disturbance
& $\omega_t$
& Target's yaw rate
& rad/s\\

Measured output
& $\hat{\rho}-d^\ast$
& Estimated range deviation obtained from the acoustic measurements
& m\\

Measured output
& $\hat{\beta}$
& Estimated relative bearing obtained from the acoustic measurements
& rad\\
\hline
\end{tabular}
\end{table*}

The range estimate is obtained from the amplitude model rather than from a direct distance measurement (\autoref{eq:rho}). 

